\section{Task 3: Jointly Trained Soft Mixture of Experts}\label{sec:t3}

\subsection{Model and losses}
Hard routing fails when the classifier is uncertain or when two corruptions mix. The soft model (Fig.~\ref{fig:arch23}, bottom) replaces the decision by a differentiable combination. A gate $G$ with the classifier's architecture produces weights
\begin{equation}
w=\mathrm{softmax}\bigl(G(\tilde{x})/\tau\bigr),\qquad \hat{x}=w_0\tilde{x}+w_1A_{salt}(\tilde{x})+w_2A_{blur}(\tilde{x})+w_3A_{occ}(\tilde{x}),
\end{equation}
where branch~0 is the identity and $\tau$ controls whether the weights are sharp or spread out. Because the mixture is differentiable, gate and experts are trained jointly from the reconstruction error. The components are \emph{not} initialised randomly: the gate starts from the trained Task~2 classifier and the experts from the final specialists. Training has two stages: a short warm-up (2 epochs, learning rate $10^{-3}$) in which the experts are frozen, in evaluation mode, and only the gate is trained, followed by 30 epochs of joint fine-tuning with a much smaller learning rate. The loss is
\begin{equation}
\mathcal{L}_{MoE}=\lambda_1\mathcal{L}_{1}+\lambda_s(1-\mathrm{SSIM})+\lambda_c\mathcal{L}_{CE}+\lambda_b\mathcal{L}_{balance},
\end{equation}
with $\mathcal{L}_{balance}=\sum_{k=0}^{3}(\bar{w}_k-\tfrac14)^2$, where $\bar{w}_k$ is the mean weight of branch $k$ in a mini-batch. The cross-entropy acts on the raw gate logits, which keeps the question ``does the gate know the true class'' separate from ``how sharply does it mix'' (the temperature). We tie $\lambda_1=\alpha$ and $\lambda_s=1-\alpha$, so that the five quantities Optuna tuned (joint learning rate, temperature, $\lambda_c$, $\lambda_b$ and the L1/SSIM weighting $\alpha$) map one-to-one onto the assignment's list and the reconstruction weights cannot both shrink to zero. The selected values are $\tau=2.68$, $\lambda_c=0.48$, $\lambda_b=1.20$, $\alpha=0.61$ and a joint learning rate of $4.3\times10^{-5}$; median pruning removed half of the trials (Section~\ref{sec:optuna}). The first mixture (with $\alpha=0.32$) inherited the colour loss of its experts, and the final one was retrained from the corrected specialists (Section~\ref{sec:colour}).

\begin{figure*}[t]
\centering
\includegraphics[width=\textwidth]{curves_task3.pdf}
\caption{Task 3 training: validation objective and gate accuracy for the first (dotted) and final runs, and validation SSIM over the warm-up and joint stages of the final run.}\label{fig:curves3}
\end{figure*}

\subsection{Reconstruction results}
Table~\ref{tab:task3} compares the three restoration systems on every test condition. The soft mixture reaches an overall SSIM of \nSoftSSIM{} against \nHardPredSSIM{} for predicted hard routing and \nTOneSSIMall{} for the universal model. Clean images stay almost untouched (SSIM \nSoftSSIMclean{}) because the gate puts most of its weight on the identity branch. On the corrupted types the mixture is on par with hard routing and with the universal model (overall \nSoftSSIMsaltpepper{} for noise, \nSoftSSIMblur{} for blur and \nSoftSSIMocclusion{} for occlusion, PSNR \nSoftPSNRsaltpepper{}, \nSoftPSNRblur{} and \nSoftPSNRocclusion{}\,dB); it is slightly better than hard routing at the low severities (for example blur $k=3$: 0.820 against 0.792) and slightly worse at the high ones (noise $p=0.15$: 0.778 against 0.790), which is what a smoother, less committed decision would be expected to do. Blending experts does not remove the fundamental limit on heavy occlusion.

\begin{table*}[t]
\centering
\caption{Soft mixture of experts (Task~3) compared with the universal model and predicted hard routing on the full test protocol.}\label{tab:task3}
\small
\adjustbox{max width=\textwidth}{\begin{tabular}{llcccrr}
\toprule
 & & \multicolumn{3}{c}{SSIM} & \multicolumn{2}{c}{soft MoE} \\
\cmidrule(lr){3-5}\cmidrule(lr){6-7}
Corruption & Level & universal & hard (pred.) & soft MoE & L1 & PSNR (dB) \\
\midrule
clean & -- & 0.814 & 0.996 & 0.991$\pm$0.017 & 0.0081 & 40.19 \\
salt-and-pepper & low & 0.813 & 0.796 & 0.807$\pm$0.070 & 0.0427 & 24.98 \\
salt-and-pepper & medium & 0.808 & 0.794 & 0.795$\pm$0.075 & 0.0436 & 24.83 \\
salt-and-pepper & high & 0.796 & 0.790 & 0.778$\pm$0.075 & 0.0441 & 24.71 \\
Gaussian blur & low & 0.814 & 0.792 & 0.820$\pm$0.074 & 0.0400 & 25.55 \\
Gaussian blur & medium & 0.805 & 0.790 & 0.795$\pm$0.086 & 0.0424 & 25.04 \\
Gaussian blur & high & 0.751 & 0.757 & 0.754$\pm$0.096 & 0.0448 & 24.49 \\
occlusion & low & 0.768 & 0.747 & 0.771$\pm$0.083 & 0.0484 & 22.78 \\
occlusion & medium & 0.718 & 0.699 & 0.699$\pm$0.093 & 0.0575 & 21.28 \\
occlusion & high & 0.634 & 0.625 & 0.621$\pm$0.098 & 0.0718 & 19.16 \\
\midrule
all (36,690) & & 0.772 & 0.779 & 0.783$\pm$0.120 & 0.0443 & 25.30 \\
\bottomrule
\end{tabular}
}
\end{table*}

\subsection{Behaviour of the gating network}
\textbf{Average weights.} Table~\ref{tab:gate} and Fig.~\ref{fig:heat} give the mean routing weight of each branch for every true corruption and severity. The gate is well aligned with the truth: the matching expert receives on average \nHomeWeightSalt{} of the weight for salt-and-pepper, \nHomeWeightBlur{} for blur and \nHomeWeightOcc{} for occlusion, and the identity branch \nHomeWeightClean{} for clean images. The weight grows with severity (occlusion: from \nGateOccLow{} at 10\% coverage to \nGateOccHigh{} at 35\%), as it should, since heavier corruption is easier to recognise. Weights are not one-hot because the temperature (2.68) deliberately spreads them: the mean largest weight is \nGateMaxMean{}.

\begin{table}[t]
\centering
\caption{Mean gate weights by true condition (full test protocol).}\label{tab:gate}
\small
\adjustbox{max width=\columnwidth}{\begin{tabular}{lcccc l}
\toprule
True condition & $w_0$ identity & $w_1$ salt & $w_2$ blur & $w_3$ occl. & dominant \\
\midrule
clean  & 0.808 & 0.030 & 0.098 & 0.063 & clean \\
salt-and-pepper low & 0.052 & 0.838 & 0.076 & 0.034 & salt-and-pepper \\
salt-and-pepper medium & 0.041 & 0.865 & 0.068 & 0.026 & salt-and-pepper \\
salt-and-pepper high & 0.042 & 0.864 & 0.067 & 0.027 & salt-and-pepper \\
Gaussian blur low & 0.148 & 0.017 & 0.807 & 0.028 & Gaussian blur \\
Gaussian blur medium & 0.096 & 0.014 & 0.865 & 0.025 & Gaussian blur \\
Gaussian blur high & 0.096 & 0.014 & 0.864 & 0.025 & Gaussian blur \\
occlusion low & 0.108 & 0.064 & 0.028 & 0.800 & occlusion \\
occlusion medium & 0.012 & 0.026 & 0.005 & 0.958 & occlusion \\
occlusion high & 0.002 & 0.009 & 0.001 & 0.988 & occlusion \\
\bottomrule
\end{tabular}
}
\end{table}

\begin{figure}[t]
\centering
\includegraphics[width=\columnwidth]{gate_heatmap.pdf}
\caption{Routing heat-map: mean weight of each branch for each true corruption and severity.}\label{fig:heat}
\end{figure}

\textbf{Dominant versus distributed routing.} Figure~\ref{fig:soft} contrasts inputs on which one expert dominates (left: the most ``one-expert'' input of each corruption) with inputs on which the weights are spread (right). The spread-out cases are low-severity corruptions, for which the image is close to a clean one and the identity and expert branches are both plausible: the mean weight entropy is \nEntOccLow{} nats for 10\% occlusions and \nEntClean{} for clean images, against only \nEntOccHigh{} for 35\% occlusions, where the gate is nearly one-hot (Table~\ref{tab:gate}). For such inputs the mixture yields a smoother output than a hard decision would.

\begin{figure*}[t]
\centering
\includegraphics[width=0.94\textwidth]{soft_weights_examples.pdf}
\caption{Left: inputs for which a single expert dominates the mixture. Right: inputs for which the weights are distributed over several branches.}\label{fig:soft}
\end{figure*}

\textbf{Inactive or dominant experts.} Table~\ref{tab:dominance} lists, for each true type, the percentage of inputs on which each branch has the largest weight. No branch is dead: each branch is the top branch for a substantial share of inputs (overall: identity \nTopShareclean{}\%, salt-and-pepper \nTopSharesaltpepper{}\%, blur \nTopShareblur{}\%, occlusion \nTopShareocclusion{}\%), and no branch dominates unrelated inputs, so the balance regulariser did its job without forcing a uniform mixture.

\begin{table}[t]
\centering
\caption{Dominance analysis: where each branch has the largest weight.}\label{tab:dominance}
\small
\adjustbox{max width=\columnwidth}{\begin{tabular}{lccccc}
\toprule
True type & \multicolumn{4}{c}{\% of inputs where branch has the largest weight} & mean max weight \\
\cmidrule(lr){2-5}
 & identity & salt & blur & occl. & \\
\midrule
clean & 99.1 & 0.0 & 0.7 & 0.2 & 0.811 \\
salt-and-pepper & 0.0 & 100.0 & 0.0 & 0.0 & 0.856 \\
Gaussian blur & 0.0 & 0.0 & 100.0 & 0.0 & 0.846 \\
occlusion & 1.9 & 0.0 & 0.0 & 98.0 & 0.921 \\
\midrule
all inputs & 10.5 & 30.0 & 30.1 & 29.4 & 0.868 \\
\bottomrule
\end{tabular}
}
\end{table}

\begin{figure*}[t]
\centering
\includegraphics[width=0.94\textwidth]{examples_soft.pdf}
\caption{Twelve representative test inputs for the soft mixture of experts.}\label{fig:ex3}
\end{figure*}

\begin{figure}[t]
\centering
\includegraphics[width=\columnwidth]{failure_soft.pdf}
\caption{The four lowest-SSIM corrupted test inputs for the soft mixture.}\label{fig:fail3}
\end{figure}

\textbf{Failure cases.} Figure~\ref{fig:fail3} shows the four worst outputs of the mixture (SSIM 0.24--0.31). They are the same kind of input as for the other systems: 35\% occlusions that hide the subject. The first is the cat on a black background whose black mask merges with the background, so the input looks like a half-empty image and the mixture produces a grey-brown cloud; in the other three (a beagle, a basset hound and a dark terrier on grass) the gate gives the occlusion expert nearly all of the weight, which is the correct routing, but the expert cannot recreate the animal. The worst cases of the mixture are about 0.03--0.04 SSIM below those of the universal model (0.24--0.31 against 0.27--0.35). Because the information is absent, no routing or blending strategy can recover these inputs, and all three systems produce a smooth, plausible fill instead of the subject.
